Over the past decade, Artificial Intelligence (AI) has been firmly established in our society widely adopted across various sectors by optimizing operations, supporting decision-making processes, and enhancing productivity through the reduction of operational inefficiencies [1] \cite{Ali2023}. In healthcare, AI supports disease detection and diagnosis, and [2] \cite{PanForoughi2024} enables faster and more accurate medical interventions. In finance, AI is employed to detect fraud and manage customized financial services [3]\cite{OyeniyiEtAl2024}, decide eligibility in loan approval. Law enforcement agencies utilize artificial intelligence to aid in decision making and crime prevention. [4] \cite{SitumeangEtAl2024}, [5] \cite{JatnaManthovaniHasbullah2024}, while the automotive sector leverages AI for the development and operation of autonomous vehicles [6]\cite{10495655}. This widespread adoption is underscored by recent industry reports. According to McKinsey’s “The state of AI in 2025: Agents, innovation, and transformation”, 88\% organizations now use AI in at least one business function, with generative AI important focus for many companies. The Artificial Intelligence Index Report indicates that, from 2019 to 2024 worldwide, the number of funded AI companies increased by 60\%, meaning robust investment and innovation in the field. [7]\cite{McKinseySuperagency2025} Moreover, when observing the use of AI within the companies reveal more than threefold increase in the number of employees are adopting generative AI into their works when comparing with the previous year's [8] \cite{MaslejEtAl2025}Despite the substantial advancements and pervasive integration of AI, experts in this field are increasingly addressing critical questions : Should we trust blindly in these AI-generated predictions and results? Does AI support human-centered approaches? Can we evaluate the explainability and transparency of AI? Such inquiries highlighted, the inherent need for robust methodologies and frameworks for establishing AI ethics.[9] \cite{DeSilvaHalloluwaVyas2025}. The need is crucial to overcome the problem of opacity in AI results, referred to as “black box”. A powerful technique is deep learning, which uses multi-layered neural networks to process complex data, however, the intricate nature of these networks often causes their decision-making process to be opaque, which hinders the establishment of trust in AI systems. This opaqueness becomes particularly problematic in the fields of healthcare, finance, and autonomous systems, where decisions have a significant impact on human life.

Consequently, the demand for Explainable AI has urged us to enhance the transparency and interpretability inherent in complex AI models [10] \cite{Yang2023SurveyOE}. Explainable Artificial Intelligence (XAI) refers to the set of methods and techniques that aim to present AI systems' decisions and functionalities in a human-understandable manner. Specifically, in image classification tasks, the use of Explainable AI techniques has become the primary focus through auditing model behavior. Post-hoc approaches  in XAI focus on methods to simplify decision making logic, that are not intrinsically understandable to human users.[12] \cite{Zschech2025Inherently} These techniques, most notably Saliency Maps and Attribution Methods, aim to focus attention on the specific pixels or spatial regions that triggered a classification such as identifying the texture of a lesion in a medical scan, spatial sequence in biopsy resolution which allows to identify the risk factor in pregnancy or the facial recognition from images or video footage, or the shape of a pedestrian in autonomous driving.

However, despite their visual appeal, there is a growing consensus that many existing XAI methods fail to meet expectations, generally due to the technical limitations or lack of human-centered design.[11] \cite{bove2025explanationsfailtypologydiscussion} A preliminary study conducted using the CIFAR-10 dataset and the LIME (Local Interpretable Model-agnostic Explanations) explainer, anomalous explanation patterns were observed. Precisely, the explainer highlighted obscured local areas, where the areas are not relevant features of the objects themselves. This anomaly influenced class predictions, suggesting that the model was not performing generalized feature extraction but memorizing specific, irrelevant pixel-level artifacts from the training data.

Such observations address the lack of underlying memorization phenomena , where models utilize superficial correlations in the training data instead of learning robust representations. This behavior stress concerns about the faithfulness of post-hoc explanations, and their ability reflect model reasoning.

Consequently, this thesis seeks a bridge this gap by investigating the extent to which model memorization results in a misleading post-hoc explanations, using these observed anomalies as a baseline, extending this baseline with different post-hoc explanation methods to develop more faithful auditing frameworks for image classification models.

Figure 1 shows the overall structure of this thesis.

This thesis is organized as follows:

– Chapter 1: Introduction introduces the fundamental questions regarding the necessity of Explainable AI (XAI), the reliability of these systems in critical sectors, and the specific impact of model memorization on the fidelity of XAI outputs.

– Chapter 2: Literature Review and Background After a deep understanding and research on the context, ‘Literature Review and Background’ section proceed in detail with relevant articles related Image Classification, Explainable AI, Memorization Phenomenon.

– Chapter 3: Methodology section, provides detailed mathematical approach of explainable ai techniques, and CNN architecture that will be used in Implementation

– Chapter 4: Implementation and Structure section, provides the implementation process of the thesis which will introduced in appendix.

– Chapter 5: Evaluation and Results section evaluates the implementation results.

– Chapter 6: Conclusion and Future Work Finally, the whole project is concluded, and the future works are provided in section ‘Conclusion and Future Work’.
